% Irodov Part VI §6.2 original problems 81--96, closing §6.2.
% Original printed pp.241--243, user PDF pages 250--252;
% answer pp.342--344, user PDF pages 351--353.
% Scattering and barrier diagram conventions read from original figures 6.2--6.6.
\wavequantproblem{81}{241}
质量为$m$的粒子被局限在边长为$l_1,l_2$的二维
无限深矩形势阱中。(a) 求全部允许能量；
(b) 若势阱为边长$l$的正方形，求前四个不同能级。
\exsolution
势阱内的定态薛定谔方程可分离：
$\psi(x,y)=X(x)Y(y)$。
边界上波函数为零，故
$X\propto\sin(n_1\pi x/l_1)$、
$Y\propto\sin(n_2\pi y/l_2)$，
$n_1,n_2=1,2,\ldots$。得到
\[
 \boxed{E_{n_1n_2}=\frac{\pi^2\hbar^2}{2m}
 \left(\frac{n_1^2}{l_1^2}+\frac{n_2^2}{l_2^2}\right)}.
\]
正方形时令
$\varepsilon=\pi^2\hbar^2/(2ml^2)$；
最低四个\emph{不同}的$n_1^2+n_2^2$为
$2,5,8,10$，因此
\[
 \boxed{E_1=2\varepsilon,\quad E_2=5\varepsilon,\quad
 E_3=8\varepsilon,\quad E_4=10\varepsilon}.
\]
\excheck 二维正方阱中$(n_1,n_2)$与$(n_2,n_1)$
常产生简并；“第四能级”应按不同能量而非重复状态排序。

\wavequantproblem{82}{241}
粒子在二维无限深矩形阱
$0<x<a,\ 0<y<b$中处于基态，
求其位于$0<x<a/3$区域的概率。
\exsolution
归一化基态波函数
$\psi_{11}=\sqrt{4/(ab)}
\sin(\pi x/a)\sin(\pi y/b)$。
对全部$y$积分为1，
故
\[
 P=\frac2a\int_0^{a/3}\sin^2\frac{\pi x}{a}\,dx
 =\boxed{\frac13-\frac{\sqrt3}{4\pi}
 \simeq0.1955}.
\]
\excheck 由于基态在$x=0$的概率密度为0，
左侧三分之一的概率比几何面积份额$1/3$小。

\wavequantproblem{83}{241}
质量$m$的粒子处于边长$a$的三维无限深立方势阱。
求：(a)全部本征能量；
(b)第三、第四不同能级的间隔；
(c)第六个不同能级的能量及简并度。
\exsolution
沿三坐标可分离并施加端点节点条件，
\[
 \boxed{E_{n_1n_2n_3}=
 \frac{\pi^2\hbar^2}{2ma^2}
 (n_1^2+n_2^2+n_3^2)},\quad n_j\ge1.
\]
最低六个\emph{不同}的平方和依次为
$3,6,9,11,12,14$，例如第三能级属于$(1,2,2)$，
第四能级属于$(1,1,3)$。
记$\varepsilon=\pi^2\hbar^2/(2ma^2)$，
所以
\[
 \boxed{E_4-E_3=(11-9)\varepsilon
 =\frac{\pi^2\hbar^2}{ma^2}}.
\]
第六能级对应$(1,2,3)$的六种排列，
\[
 \boxed{E_6=14\varepsilon=\frac{7\pi^2\hbar^2}{ma^2}},
 \qquad \boxed{g_6=3!=6}.
\]
\excheck 简并度不含自旋；
若额外计自旋自由度，需要另外乘相应因子。

\wavequantproblem{84}{241}
从定态薛定谔方程证明：
粒子势能$U(x)$在某点发生\emph{有限跃变}时，
波函数一阶导数仍连续。
\exsolution
定态方程
\[
 -\frac{\hbar^2}{2m}\psi''(x)+U(x)\psi(x)=E\psi(x).
\]
将方程在跃变点$x_0$两侧
$[x_0-\delta,x_0+\delta]$上积分，
\[
 \psi'(x_0+\delta)-\psi'(x_0-\delta)
 =\frac{2m}{\hbar^2}
 \int_{x_0-\delta}^{x_0+\delta}(U-E)\psi\,dx.
\]
若$U$有界且$\psi$有限，
右边随$\delta\to0$趋于零，
因此
\[
 \boxed{\psi'(x_0^+)=\psi'(x_0^-)}.
\]
\excheck 有限跃变也要求$\psi$连续；
$\delta$型奇异势不属于本题，
其$\psi'$可出现有限跳跃。

\wavequantproblem{85}{241}
粒子位于图6.2所示的单侧有限深方势阱：
$x\le0$处为无限高势壁，
$0<x<l$有$U=0$，
$x\ge l$有$U=U_0>0$。
(a) 推导$E<U_0$的束缚能级方程并表明能谱离散；
(b) 求存在第$n$个束缚态所需$l^2U_0$的阈值。
\exsolution
令
$k=\sqrt{2mE}/\hbar$、
$\kappa=\sqrt{2m(U_0-E)}/\hbar$。
归一化束缚解在阱内、阱外分别是
$\psi_<=A\sin(kx)$、
$\psi_>=B e^{-\kappa(x-l)}$。
在$x=l$令波函数及一阶导数连续，
\[
 A\sin(kl)=B,\quad
 Ak\cos(kl)=-B\kappa,
\]
于是
\[
 \boxed{k\cot(kl)=-\kappa
 \quad\Longleftrightarrow\quad \tan(kl)=-\frac k\kappa}.
\]
又$k^2+\kappa^2=2mU_0/\hbar^2$，
\[
 \boxed{\sin(kl)=\pm\frac{\hbar k}{\sqrt{2mU_0}}},
\]
同时必须选取$\tan(kl)<0$的根。
离散实根给出离散本征能级。
在第$n$个态刚好出现时$E\uparrow U_0$、
$\kappa\downarrow0$，故
$kl=(n-\tfrac12)\pi$，从而
\[
 \boxed{(l^2U_0)_{n,\mathrm{th}}=
 \frac{(2n-1)^2\pi^2\hbar^2}{8m}},
 \quad n=1,2,\dots.
\]
\excheck 阈值处$E=U_0$的态已不严格平方可积，
因此真正束缚态须满足严格大于阈值。
高斯/文本抽取常漏去$(2n-1)^2$的平方。

\wavequantproblem{86}{241}
沿用前题单侧有限深阱。
当$E=U_0/2$且
$l^2U_0=(3\pi/4)^2\hbar^2/m$时，
求粒子出现在$x>l$区域的概率。
\exsolution
此时
$k=\kappa=\sqrt{mU_0}/\hbar$，
$kl=3\pi/4$，
满足$\tan(kl)=-1$。
取
\[
 \psi(x)=
 \begin{cases}
 A\sin(kx),&0<x<l,\\
 A\sin(kl)e^{-k(x-l)},&x>l.
 \end{cases}
\]
阱外积分
$I_{\rm out}=|A|^2\sin^2(kl)/(2k)=|A|^2/(4k)$；
阱内积分
\[
 I_{\rm in}=|A|^2
 \left[\frac l2-\frac{\sin(2kl)}{4k}\right]
 =\frac{|A|^2}{4k}\left(1+\frac{3\pi}{2}\right).
\]
所以
\[
 \boxed{P_{\rm out}=
 \frac{I_{\rm out}}{I_{\rm out}+I_{\rm in}}
 =\frac2{4+3\pi}\simeq0.149}.
\]
\excheck 类禁区概率不为零是波函数指数尾部的结果，
不是经典粒子穿透并持续向右运动的定态概率流。

\wavequantproblem{87}{242}
质量$m$的粒子在半径$r_0$的三维球形
无限深势阱中运动：
$U=0$（$r<r_0$）、
$U=\infty$（$r\ge r_0$）。
若波函数仅与径向$r$有关，求允许能量。
\exsolution
径向定态方程（角动量$l_{\rm orb}=0$）为
\[
 \psi''+\frac2r\psi'+k^2\psi=0,\quad
 k^2=\frac{2mE}{\hbar^2}.
\]
令$\chi(r)=r\psi(r)$，则
$\chi''+k^2\chi=0$。
原点波函数有限要求$\chi(0)=0$，
球壁要求$\chi(r_0)=0$，
因此
$\chi\propto\sin(kr)$且$kr_0=n\pi$。
所以
\[
 \boxed{E_n=\frac{n^2\pi^2\hbar^2}{2mr_0^2}},
 \quad n=1,2,\dots.
\]
\excheck 这是只考虑球对称$s$波的能谱，
并非所有角动量分波的完整三维球阱谱。

\wavequantproblem{88}{242}
在前题球形无限深势阱中：
(a) 求只依赖$r$的归一化波函数；
(b) 求基态最概然径向$r_{\rm mp}$以及
$P(r<r_{\rm mp})$。
\exsolution
设$\psi_n=A_n\sin(n\pi r/r_0)/r$。
归一化积分
\[
 1=4\pi|A_n|^2\int_0^{r_0}
 \sin^2\frac{n\pi r}{r_0}\,dr
 =2\pi r_0|A_n|^2,
\]
故
\[
 \boxed{\psi_n(r)=\frac1{\sqrt{2\pi r_0}}
 \frac{\sin(n\pi r/r_0)}r},\quad r<r_0.
\]
基态的径向概率密度
$dP/dr=4\pi r^2|\psi_1|^2
=(2/r_0)\sin^2(\pi r/r_0)$，
在$r=r_0/2$达最大。因此
\[
 \boxed{r_{\rm mp}=\frac{r_0}2},\qquad
 \boxed{P(r<r_{\rm mp})=\frac2{r_0}
 \int_0^{r_0/2}\sin^2\frac{\pi r}{r_0}dr=\frac12}.
\]
\excheck 最概然\emph{径向}应最大化$4\pi r^2|\psi|^2$，
不是单独最大化$|\psi(r)|^2$。

\wavequantproblem{89}{242}
粒子在球形有限深势阱中：
$U(r)=0$（$r<r_0$）、
$U(r)=U_0$（$r\ge r_0$）。
(a) 对球对称定态导出$E<U_0$束缚态能级条件；
(b) 求出现第一个束缚态所需$r_0^2U_0$阈值。
\exsolution
仍令$\chi=r\psi$，
则径向方程化为一维形式。
阱内
$\chi_<=A\sin(kr)$，
阱外取衰减解
$\chi_>=B e^{-\kappa(r-r_0)}$，
其中$k=\sqrt{2mE}/\hbar$、
$\kappa=\sqrt{2m(U_0-E)}/\hbar$。
在$r=r_0$匹配$\chi,\chi'$，
\[
 \boxed{k\cot(kr_0)=-\kappa},
 \quad
 \boxed{\sin(kr_0)=\pm\frac{\hbar k}{\sqrt{2mU_0}}}.
\]
第一次出现束缚态的阈值在
$\kappa\to0$、$kr_0\to\pi/2$，
\[
 \boxed{(r_0^2U_0)_{\mathrm{th}}=
 \frac{\pi^2\hbar^2}{8m}}.
\]
\excheck 对三维\emph{球形}有限势阱必须超过一定深度
才支持$s$波束缚态；
这一点不同于对称的一维有限势阱。

\wavequantproblem{90}{242}
质量为$m$的粒子处于一维简谐势
$U(x)=kx^2/2$。
已知其基态波函数形式为
$\psi(x)=A e^{-\alpha x^2}$（$\alpha>0$）。
从薛定谔方程求$\alpha$与基态能量。
\exsolution
计算导数
$\psi''=(4\alpha^2x^2-2\alpha)\psi$。
代入
\[
 -\frac{\hbar^2}{2m}\psi''+\frac{kx^2}{2}\psi=E\psi
\]
得
\[
 \left[\frac{\hbar^2\alpha}{m}
 +\left(\frac k2-\frac{2\hbar^2\alpha^2}{m}\right)x^2\right]\psi
 =E\psi.
\]
因等式对一切$x$成立，$x^2$系数为0，
$\alpha=\sqrt{mk}/(2\hbar)$，
而常数项即能量。故
\[
 \boxed{\alpha=\frac{m\omega}{2\hbar}},\quad
 \boxed{E_0=\frac{\hbar\omega}{2}},
 \qquad \omega=\sqrt{k/m}.
\]
\excheck 这是严格量子基态，比粗略不确定度估算题73
所得$\sim\hbar\omega$多出明确的$1/2$系数。

\wavequantproblem{91}{242}
氢原子电子某一球对称定态的波函数为
$\psi(r)=A(1+ar)e^{-\alpha r}$，其中
$A,a,\alpha$均为常数。求其能量及主量子数。
\exsolution
氢原子库仑势
$U=-k_ee^2/r$。
已知球对称态的径向薛定谔方程
\[
 -\frac{\hbar^2}{2m_e}
 \left(\psi''+\frac2r\psi'\right)
 -\frac{k_ee^2}{r}\psi=E\psi.
\]
令$f=(1+ar)e^{-\alpha r}$，
则
\[
 f'=e^{-\alpha r}(a-\alpha-a\alpha r),\quad
 f''=e^{-\alpha r}(\alpha^2-2a\alpha+a\alpha^2r).
\]
代入并分别比较常数、$r^{-1}$及$r$项，
求得
$\alpha=1/(2a_0)$、$a=-1/(2a_0)$，
其中$a_0=\hbar^2/(m_ek_ee^2)$。
于是该态正是氢原子的$2s$径向态，
\[
 \boxed{E=-\frac{m_e(k_ee^2)^2}{8\hbar^2}
 =-\frac{13.6}{4}\,\mathrm{eV}=-3.40\,\mathrm{eV}},
 \qquad\boxed{n=2}.
\]
\excheck 虽然波函数只有$r$，仍可有$n=2$的径向节点，
不能由角动量为0就断定$n=1$。

\wavequantproblem{92}{242}
基态氢原子电子波函数
$\psi=Ae^{-r/a_0}$，
$a_0$为第一玻尔半径。
求：(a)最概然核间距离；
(b)电子所受库仑力模的平均值；
(c)平均势能。
\exsolution
归一化给
$A=1/\sqrt{\pi a_0^3}$。
径向概率密度
$f(r)=4r^2e^{-2r/a_0}/a_0^3$，
令$f'(r)=0$得
\[
 \boxed{r_{\rm mp}=a_0}.
\]
库仑力模为$F(r)=k_ee^2/r^2$，
所以
\[
 \langle F\rangle=
 \int_0^\infty f(r)\frac{k_ee^2}{r^2}\,dr
 =\boxed{\frac{2k_ee^2}{a_0^2}}.
\]
势能$U(r)=-k_ee^2/r$，则
\[
 \boxed{\langle U\rangle=
 -\frac{k_ee^2}{a_0}}.
\]
\excheck 力模的\emph{平均值}不是
$k_ee^2/\langle r\rangle^2$；
必须先按波函数求概率平均。

\wavequantproblem{93}{242}
基态氢原子电子波函数
$\psi=Ae^{-r/a_0}$。
求在氢原子\emph{原子核中心}，
该电子电荷云产生的平均静电势。
\exsolution
电荷密度
$\rho(r)=-e|\psi(r)|^2$，
其中$|A|^2=1/(\pi a_0^3)$。
球对称分布在原点的势为
\[
 \varphi(0)=k_e\int_{\mathbb R^3}
 \frac{\rho(\boldsymbol r)}r\,d^3r
 =4\pi k_e\int_0^\infty \rho(r)\,r\,dr.
\]
代入积分
$\int_0^\infty r e^{-2r/a_0}\,dr=a_0^2/4$，
得
\[
 \boxed{\varphi(0)=-\frac{k_ee}{a_0}
 \simeq-27.2\,\mathrm V}.
\]
\excheck 此处仅计算\emph{电子云}在核位置的势，
不包括点核自身的发散自势。

\wavequantproblem{94}{243}
图6.3中势阶跃
$U(x)=0$（$x<0$），
$U(x)=U_0$（$x>0$）。
求：(a)当$E>U_0$时从左方入射粒子的反射系数；
(b)当$E<U_0$时，右侧禁区中概率密度降到界面处
$1/e$所需渗入深度。
\exsolution
(a) 两区域波数
$k_1=\sqrt{2mE}/\hbar$、
$k_2=\sqrt{2m(E-U_0)}/\hbar$。
写
$\psi_1=Ae^{ik_1x}+Be^{-ik_1x}$、
$\psi_2=Ce^{ik_2x}$。
在$x=0$匹配$\psi$与$\psi'$，
由$A+B=C$、
$k_1(A-B)=k_2C$，
\[
 \boxed{R=\frac{|B|^2}{|A|^2}
 =\left(\frac{k_1-k_2}{k_1+k_2}\right)^2}.
\]
(b) 若$E<U_0$，
$\psi_2=C e^{-\kappa x}$，
$\kappa=\sqrt{2m(U_0-E)}/\hbar$。
概率密度衰减为$e^{-2\kappa x}$，
满足$e^{-2\kappa x_{\rm eff}}=e^{-1}$，
故
\[
 \boxed{x_{\rm eff}=\frac1{2\kappa}
 =\frac{\hbar}{2\sqrt{2m(U_0-E)}}}.
\]
\excheck 本题“下降$e$倍”的对象是\emph{概率密度}，
不是波函数振幅；两者的衰减长度相差2倍。

\wavequantproblem{95}{243}
利用半经典WKB公式，求能量为$E<U_0$的电子穿过
图6.4中的矩形势垒和图6.5中的线性上升三角势垒的
透射概率指数近似。矩形宽$l$高$U_0$；
三角势垒取$U(x)=U_0(x+l)/l$于$-l<x<0$，
在$x=0$降至零。
\exsolution
WKB禁止区透射率的主指数是
\[
 D\simeq\exp\!\left[-\frac2\hbar
 \int_{x_-}^{x_+}\sqrt{2m(U(x)-E)}\,dx\right].
\]
(a) 矩形势垒禁止区贯穿整个宽度，
\[
 \boxed{D_{\rm rect}\simeq
 \exp\!\left[-\frac{2l}{\hbar}\sqrt{2m(U_0-E)}\right]}.
\]
(b) 三角势垒的左转折点
$x_t=-l+lE/U_0$。
由$dU=(U_0/l)\,dx$，
\[
 \int_{x_t}^{0}\sqrt{2m(U-E)}\,dx
 =\frac{l\sqrt{2m}}{U_0}
 \int_E^{U_0}\sqrt{U-E}\,dU
 =\frac{2l\sqrt{2m}}{3U_0}(U_0-E)^{3/2}.
\]
故按图6.5所示\emph{单侧}斜坡，
\[
 \boxed{D_{\rm tri}\simeq
 \exp\!\left[-\frac{4l\sqrt{2m}}{3\hbar U_0}
 (U_0-E)^{3/2}\right]}.
\]
\excheck 用户所附版本书末第95(b)题印有
$8l/(3\hbar U_0)$而非本式的$4l/(3\hbar U_0)$；
依原图仅一个宽度$l$的线性斜坡时，严格的WKB
积分是系数4；若将禁区积分路径重复两次才出现8。
该冲突应保留供二校，不能悄悄修改图或积分区间。

\wavequantproblem{96}{243}
用WKB公式求质量$m$、能量$E<U_0$的粒子
穿过图6.6所示抛物势垒
$U(x)=U_0(1-x^2/l^2)$的透射概率指数近似。
\exsolution
经典转折点由$U(x_\pm)=E$求出
$x_\pm=\pm a$，
$a=l\sqrt{1-E/U_0}$。
在禁止区
$U(x)-E=(U_0-E)(1-x^2/a^2)$，
因此
\[
 \int_{-a}^{a}\sqrt{2m(U-E)}\,dx
 =a\sqrt{2m(U_0-E)}
 \int_{-1}^{1}\sqrt{1-u^2}\,du.
\]
利用半圆面积积分
$\int_{-1}^{1}\sqrt{1-u^2}\,du=\pi/2$，
WKB指数为
\[
 \boxed{D\simeq
 \exp\!\left[-\frac{\pi l}{\hbar}
 \sqrt{\frac{2m}{U_0}}\,(U_0-E)\right]}.
\]
\excheck 势垒宽度不应一律取$2l$，
应当按能量$E$确定转折点$\pm a$；
所得指数与原书第96题答案相符。
